% Source: https://artofproblemsolving.com/community/c3h2289879_floor_function1303_post_why_did_i_miss_1300
% Stable identity: aops-topic-2289879-q01
\textbf{The key fact: $g(n)-g(n-1)$ is the number of positive divisors of $n$.} For $n\ge2$, the term $\left\lfloor\dfrac{n-1}{n}\right\rfloor$ is $0$, so we may write
\[g(n)-g(n-1)=\sum_{k=1}^{n}\left(\left\lfloor\dfrac nk\right\rfloor-\left\lfloor\dfrac{n-1}k\right\rfloor\right).\]
Write $n-1=qk+r$ with $0\le r<k$. If $r<k-1$, then $n=qk+(r+1)$ with $r+1<k$, so both floors equal $q$. If $r=k-1$, then $n=(q+1)k$, so $\left\lfloor\dfrac nk\right\rfloor=q+1$. Hence each term is $1$ if $k\mid n$ and $0$ otherwise, and $g(n)-g(n-1)=d(n)$, the number of positive divisors of $n$.

\textbf{(a)} $d(7)=\boxed{2}$ (the divisors $1$ and $7$).

\textbf{(b)} $160=2^5\cdot5$, so $d(160)=(5+1)(1+1)=\boxed{12}$.

\textbf{(c)} We need $d(n)=11$. If $n=p_1^{e_1}\cdots p_s^{e_s}$, then $d(n)=(e_1+1)\cdots(e_s+1)$. Since $11$ is prime, there is only one prime factor, with $e_1+1=11$. So the solutions are exactly $\boxed{n=p^{10}\text{ for a prime }p}$ (the smallest is $2^{10}=1024$).
